\section{Conjectures and open questions}\label{open:section}

The proved statements leave a small set of next steps.  We distinguish
the one prediction stated as a conjecture from six questions whose
answers are not established here.  The definitions are recalled so
that each can be read independently.  The connection to classical
conjectures remains visible in Proposition~\ref{play:normality}:
normality of the original constant is equivalent to Goldbach together
with the naturally averaged Liouville Chowla conjecture.

\subsection{A conjecture about the typical deletion time}

In the deletion experiment of Section~\ref{fin:deletion}, begin with
all primes except $7$.  At every round retain $p_i$ only when
$2(i+1)$ is the sum of two currently retained primes; the same
prime may be used twice, and all deletions happen simultaneously.
Write $\tau(p)$ for the first round in which $p$ is absent, with
$\tau(7)=0$ and $\tau(2)=\tau(3)=\tau(5)=\infty$, and put
\[
 M(x)=\max_{5<p\le x}\tau(p),
\]
where the maximum is over primes.  Proposition~\ref{fin:depthclock}
proves $M(x)\sim\log x/\log\log x$.  This is the clock of the
last survivor in a finite window.  We conjecture that it also gives
the typical deletion time.

\begin{conjecture}[Most primes share the last survivor's clock]\label{fin:typicaldepth}
Let $P_x$ be a uniformly chosen prime in $(5,x]$.
Then
\[
 \frac{\tau(P_x)}{M(x)}\longrightarrow1
 \qquad\text{in probability}.
\]
\end{conjecture}

Here the randomness is only the uniform choice of $P_x$; the deletion
process itself is deterministic.  The exact computation through
$10^6$ found that about $94.1\%$ of the $78\,495$ primes above $5$
disappear in rounds $11$, $12$ or $13$, with $M(10^6)=13$.
The proposition proves the clock for the last surviving prime.
The conjecture asks whether nearly the whole population follows it.
The computation is a finite experiment, not evidence for a specific
error term; fixed-depth abundance alone does not settle this question.

\subsection{How much data does arithmetic recovery need?}

Both questions in this subsection use the ordinary-integer experiment
of \eqref{rec:finite-model}.  There are $N$ labelled boxes, with hidden
contents $X_1,\ldots,X_N$ independently uniform on
$\{1,\ldots,N^2\}$.  For each unordered pair we observe one bit saying
whether the contents are coprime, independently flipped with common
probability $\varepsilon_N$.  The flips are independent of the
integers; the decoder is not told $\varepsilon_N$.
Theorem~\ref{rec:threshold} identifies the consistency threshold
\[
 T_N=N(1-2\varepsilon_N)^2\longrightarrow\infty.
\]
It names the shared primes on a fraction of pairs tending to one.
Corollary~\ref{rec:profiles} then recovers divisibility by each fixed
prime at a fixed labelled box.  These are asymptotic results; the
following asks for a useful finite guarantee.  Take $N^2\ge p$ so
that conditioning on either divisibility outcome is meaningful.

\begin{question}[How many boxes are needed to name a prime?]\label{rec:prime-sample-question}
In the finite experiment, fix a prime $p$, a constant edge-flip rate
$\varepsilon_N=\varepsilon\ne1/2$, and
$0<\delta<1/2$. How large must $N$ be to decide whether $p\mid X_i$,
with conditional error at most $\delta$ both for a multiple and for a
nonmultiple of $p$? Require a computationally feasible procedure that
is not supplied the error rate. How does the answer grow with $p$
and with the closeness of the channel to a fair coin?
\end{question}

The second question concerns the original message rather than the
full collection of shared-prime sets.  Fix an unknown mask
$g=(g_j)$ with $g_1=1$, and set
\[
 x_g(1)=0,\qquad
 x_g(n)=\frac{1-\lambda(n)}2g_{\ell(n)}\quad(n\ge2),
\]
where $\ell(n)$ is the index of the least prime factor of $n$.
Alongside the noisy graph, observe aligned digit marks
\[
 Z_i=x_g(X_i)\mathbin{\oplus}B_i',
 \qquad B_i'\sim\operatorname{Bernoulli}(\eta_N).
\]
These digit flips are independent of one another, the integers and
the edge flips.  As in Theorem~\ref{rec:message}, the digit-noise
probability $\eta_N$ is supplied; the edge-noise probability is not.
Writing $a_N=1-2\varepsilon_N$ and $b_N=1-2\eta_N$, that theorem
recovers both the message and almost every shared-prime set when
$Na_N^2\to\infty$ and $Nb_N^2\to\infty$.
Here recovering the message means recovering each fixed finite prefix
of $g$ with probability tending to one as $N$ grows.

\begin{question}[How much graph does the message need?]\label{rec:message-only-question}
If the only required output is the arithmetic message, what conditions
on the two noise sequences suffice? In particular, can the edge signal
be weaker than the threshold needed to recover shared-prime sets?
\end{question}

This asks for a separate decoding theorem.  The joint-task result
supplies the digit-noise rate to the procedure, and makes no assertion
of adaptation to an arbitrary unknown rate approaching one half.

\subsection{Can one hear every disabled prime?}

The exact bias can confuse a single disabled lane with infinitely
many of them, as Proposition~\ref{grw:bias-collision} shows.  The
whole function distinguishes that example immediately: one function
is linear and the other is transcendental.  How much farther does
this recovery go?

For this question only, write $\Phi_{\mathcal F}$ for the normalized
entire function in \eqref{fnd:two-functions} whose disabled-prime set
is $\mathcal F$.  Its Taylor coefficients are $M_k/k!$, where $M_k$
are the conditional correlations of Section~\ref{fnd:correlations}.
The counting exponent of
$\mathcal F$ is
\[
 \limsup_{R\to\infty}
 \frac{\log\max(1,\#(\mathcal F\cap[2,R]))}{\log R}.
\]
Restrict to masks with counting exponent strictly below one, so that
their zero multisets determine the normalized entire functions by
Corollary~\ref{grw:zeros}.

\begin{question}[Hearing the mask]\label{play:mask-rigidity}
Let $\mathcal F$ and $\mathcal H$ be sets of odd primes whose
counting exponents are strictly below one.  Does
\[
 \Phi_{\mathcal F}\equiv\Phi_{\mathcal H}
 \quad\Longrightarrow\quad \mathcal F=\mathcal H?
\]
Equivalently, can two different such masks have exactly the same
complex zeros, including multiplicities?
\end{question}

There are three equivalent ways to supply the data: all conditional
correlations $M_k$, the entire function $\Phi$, or the nonzero
eigenvalues of the positive mask operator $A_g$.  Indeed,
$\Psi$ is obtained from $\Phi$ by multiplying its $k$th Taylor
coefficient by the fixed nonzero number $K(k)$, and
$\Psi(z)=\prod_j(1+\alpha_jz)$ records those eigenvalues.
Thus this is also an inverse spectral question for a restricted
arithmetic family of operators.

For finite masks the answer is yes: even the first derivative
$\Phi'(0)=S_g$ determines the mask.  An infinite mask cannot agree
with a finite one because the corresponding functions have different
degrees.  A possible failure must therefore involve two infinite
masks.  We already know that identical functions would force the
same counting exponent.  The question is whether the higher
correlations retain the individual prime locations as well.

\subsection{Does the triangular match include the last root?}

Recall
\[
 J_N(z)=\sum_{k=0}^N\binom Nk K(k)z^k,
\]
where $K(k)$ is the limiting probability that $k$ random integers
are pairwise coprime.  Its roots are negative; write them as
$-R_{N,1},\ldots,-R_{N,N}$, with $R_{N,j}>0$ and multiplicities
included.  Theorem~\ref{tri:main} shows that the empirical
distribution of
\[
 \frac{\exp(e^{-\gamma}R_{N,j})}{N\log N}
\]
approaches the Dykema--Haagerup law $\operatorname{DH}_1$, whose
support is $[0,e]$; $\gamma$ is Euler's constant.  This is a
statement about the whole histogram.

\begin{question}[Does the last root find the same edge?]\label{tri:edgequestion}
Do the largest arithmetic roots satisfy
\[
 e^{-\gamma}\max_jR_{N,j}-\log N-\log\log N\longrightarrow1,
\]
or, equivalently,
\[
 \max_j\frac{\exp(e^{-\gamma}R_{N,j})}{N\log N}
 \longrightarrow e?
\]
\end{question}

This asks whether the match reaches the edge as well as the histogram.
For the Gaussian triangular matrix, the largest squared singular value
does tend to $e$ \cite[Theorem 2]{triCheliotis}.
The empirical convergence proved here leaves room for a small number
of arithmetic outliers.  It implies the corresponding lower limit is
at least $e$, since the limiting law has mass arbitrarily close to that
edge; excluding larger outliers needs more.
The same distinction leaves convergence of transformed empirical
moments unproved.

\subsection{Locating the returns to the middle height}

Let $f$ be the entire prime density of
Theorem~\ref{fnd:entire-density}, and set
$m=f(0)=1/(2\mathfrak C_2)$.
Theorem~\ref{play:midpoint} proves that every polynomial target is
met infinitely often.  For the target $m$, the origin is the only
real solution; reflection and conjugation place roots off the axes
in quartets.  The theorem does not locate those roots.

\begin{question}[Where does the middle height return?]
\label{play:midpoint-geometry}
Are there infinitely many solutions of $f(z)=m$ on the imaginary
axis?  How fast does
\[
 N_m(R)=\#\{z\in\mathbb C:|z|\le R,\ f(z)=m\},
\]
with multiplicities counted, grow as $R\to\infty$?
\end{question}

The first-quadrant calculation in Section~\ref{play:section} is an
experiment with finite approximations, not a certification of a root.
The imaginary-axis question asks about a different part of the
geometry, while the counting question measures how often the middle
height returns in the whole plane.

\subsection{The dimensions inside a singular staircase}

Let $p_j$ be the $j$th prime.  Disable exactly the lanes in
$\mathcal F_*=\{p_{2^j}:j\ge2\}$, the explicit mask
\eqref{samp:explicit-mask}, and write $\nu_*$ for its primitive
sampling measure from Theorem~\ref{samp:staircase}.  This measure is
singular continuous, has support $[0,1]$, and has Hausdorff dimension
one.  The local dimension describes a different feature: how quickly
the probability around an individual point shrinks.

For $y\in[0,1]$, let $I_n(y)$ be its dyadic interval of length
$2^{-n}$; endpoints may be assigned consistently to either adjacent
interval since the measure is nonatomic.  Choose $Y$ with law
$\nu_*$.  Then the following local dimension is itself a random
variable.

\begin{question}[The dimensions inside the staircase]
\label{play:local-dimensions}
For the explicit infinite mask
$\mathcal F_* =\{p_{2^j}:j\ge2\}$, determine the distribution of
\[
 \underline d_*(Y)=\liminf_{n\to\infty}
       \frac{-\log_2\nu_*(I_n(Y))}{n},\qquad Y\sim\nu_*.
\]
In particular, does this distribution have any atoms?
\end{question}

The example with only the lane of $3$ disabled, in
Section~\ref{play:section}, gives an exact comparison: its local
dimension equals $1$, $5/6$ or $2/3$, with probabilities $1/4$,
$1/2$ and $1/4$.  Infinitely many prime lanes may produce a richer
distribution than these three values.  The question concerns the
geometry of the construction independently of which Goldbach
answers occur.
